\subsection{Cone and cylinder of a morphism of complexes}

Let \(u:(\mathbf{C}^{\prime},d^{\prime})\to(\mathbf{C},d)\) be a morphism of complexes. Let \(\operatorname{Cyl}(u)\) and \(\operatorname{Con}(u)\) be the \(\text{graded A-modules }\operatorname{Cyl}(u)=\mathbf{C}^{\prime}\oplus\mathbf{C}^{\prime}(-1)\oplus\mathbf{C}\) , \(\operatorname{Con}(u)=\mathbf{C}^{\prime}(-1)\oplus\mathbf{C}\) and let us define graded A-linear maps of degree \((-1)\)

\(\overline{\mathbf{D}}:\operatorname {Cyl}(u)\to \operatorname {Cyl}(u),\quad \overline{\mathbf{D}} (x^{\prime},y^{\prime},x) = (d^{\prime}x^{\prime} + y^{\prime}, - d^{\prime}y^{\prime},dx - u(y^{\prime})),\)

D : Con \((u) \to \text{Con}(u)\) , D \((y', x) = (-d'y', dx - u(y'))\) .

(Here, and in what follows, we denote by \(x\) , \(y\) , \ldots{} arbitrary elements of C, \(x'\) , \(y'\) , \ldots{} arbitrary elements of C'.)

Lemma 2. --- (Cyl \((u)\) , \(\overline{\mathbf{D}}\) ) and (Con \((u)\) , D) are complexes of A-modules.

Indeed, we have

\[
\begin{array}{r l} \overline {{\mathbf {D}}} \circ \overline {{\mathbf {D}}} (x ^ {\prime}, y ^ {\prime}, x) & = \overline {{\mathbf {D}}} (d ^ {\prime} x ^ {\prime} + y ^ {\prime}, - d ^ {\prime} y ^ {\prime}, d x - u (y ^ {\prime})) = \\ & = (d ^ {\prime} (d ^ {\prime} x ^ {\prime} + y ^ {\prime}) - d ^ {\prime} y ^ {\prime}, - d ^ {\prime} (- d ^ {\prime} y ^ {\prime}), d (d x - u (y ^ {\prime})) - u (- d ^ {\prime} y ^ {\prime})) = 0 \end{array}
\]

since \(d' \circ d' = 0\) , \(d \circ d = 0\) and \(d \circ u = u \circ d'\) . Likewise \(D \circ D = 0\) .

DEFINITION 7. --- The complexes Cyl (u) and Con (u) are called respectively the cylinder and the cone of the morphism u.

Example. --- Let \(u: \mathbf{M} \to \mathbf{N}\) be a homomorphism of A-modules; then the only nonzero homogeneous components of \(\operatorname{Cyl}(u)\) and \(\operatorname{Con}(u)\) are

\[
\operatorname{Cyl} _ {1} (u) = \mathrm{M}, \quad \operatorname{Cyl} _ {0} (u) = \mathrm{M} \oplus \mathrm{N},
\]

\[
\operatorname{Con} _ {1} (u) = \mathrm{M}, \quad \operatorname{Con} _ {0} (u) = \mathrm{N},
\]

and we have \(\overline{\mathrm{D}}(m) = (m, -u(m))\) , \(\mathrm{D}(m) = -u(m)\) for \(m \in \mathbf{M}\) ; consequently,

\[
\mathrm{H} _ {1} (\text { Con } (u)) = \operatorname{Ker} (u), \quad \mathrm{H} _ {0} (\text { Con } (u)) = \operatorname{Coker} (u).
\]

The complex \(\Gamma(s)\) associated with the space \(\{s\}\) endowed with its unique cellular decomposition is reduced to the module A and is identified with a subcomplex of \(\Gamma(\text{Con}(f))\) ; denote by \(\Gamma(\text{Con}(f), s)\) the quotient complex. The map \(f\) defines a morphism of complexes \(\Gamma(f): \Gamma(X) \to \Gamma(Y)\) , and one can show that the complexes \(\Gamma(\text{Cyl}(f))\) and \(\Gamma(\text{Con}(f), s)\) are respectively identified with \(\text{Cyl}(\Gamma(f))\) and \(\text{Con}(\Gamma(f))\) .

Moreover, and without any hypotheses on X and Y, one associates to \(f\) a morphism of complexes \(f_{*}: \mathrm{C}(\mathrm{X}, \mathrm{A}) \to \mathrm{C}(\mathrm{Y}, \mathrm{A})\) . One can construct injective homotopisms from \(\mathrm{Cyl}(f_{*})\) to \(\mathrm{C}(\mathrm{Cyl}(f), \mathrm{A})\) and from \(\mathrm{Con}(f_{*})\) to \(\mathrm{C}(\mathrm{Con}(f), \{s\}, \mathrm{A})\) .

Let \(\tilde{f}: X \to \text{Cyl}(f)\) be the map which to \(x\) associates the image of \((0, x)\) , \(\alpha: Y \to \text{Cyl}(f)\) the canonical map and \(\beta: \text{Cyl}(f) \to Y\) the map which associates \(y\) to its image in \(\text{Cyl}(f)\) for \(y \in Y\) and \(f(x)\) to the image of \((t, x)\) in \(\text{Cyl}(f)\) for \(t \in [0, 1]\) and \(x \in X\) . The map \(\tilde{f}\) is a homeomorphism from \(X\) onto a closed subset of \(\text{Cyl}(f)\) , we have \(\beta \circ \alpha = \text{Id}_Y\) , and \(\alpha \circ \beta\) is topologically homotopic to the identity of \(\text{Cyl}(f)\) . These properties are to be compared with Proposition 7 below.

Consider now the graded A-linear maps of degree 0

\[
\tilde {u}: \mathrm{C} ^ {\prime} \rightarrow \operatorname{Cyl} (u), \quad \tilde {u} (x ^ {\prime}) = (x ^ {\prime}, 0, 0),
\]

\[
\alpha : \mathrm{C} \rightarrow \operatorname{Cyl} (u), \quad \alpha (x) = (0, 0, x),
\]

\[
\beta : \operatorname{Cyl} (u) \rightarrow \mathrm{C}, \quad \beta \left(x ^ {\prime}, y ^ {\prime}, x\right) = u \left(x ^ {\prime}\right) + x,
\]

\[
\pi : \mathbf {C} \rightarrow \operatorname{Con} (u), \quad \pi (x) = (0, x),
\]

\[
\tilde {\pi}: \operatorname{Cyl} (u) \to \operatorname{Con} (u), \quad \tilde {\pi} (x ^ {\prime}, y ^ {\prime}, x) = (y ^ {\prime}, x),
\]

\[
\delta : \operatorname{Con} (u) \to \mathbf {C} ^ {\prime} (- 1), \qquad \delta (y ^ {\prime}, x) = y ^ {\prime}.
\]

PROPOSITION 7. --- a) The maps \(\tilde{u}, \alpha, \beta, \pi, \tilde{\pi}, \delta\) are morphisms of complexes: we have \(u = \beta \circ \tilde{u}, \pi = \tilde{\pi} \circ \alpha, \beta \circ \alpha = 1_{\mathbb{C}}\) .

\begin{enumerate}
\def\labelenumi{\alph{enumi})}
\setcounter{enumi}{1}
\tightlist
\item
  The sequences of morphisms of complexes
\end{enumerate}

\begin{enumerate}
\def\labelenumi{(\arabic{enumi})}
\setcounter{enumi}{5}
\tightlist
\item
\end{enumerate}

\[
0 \rightarrow \mathrm{C} ^ {\prime} \stackrel {\tilde {u}} {\rightarrow} \operatorname{Cyl} (u) \stackrel {\tilde {\pi}} {\rightarrow} \operatorname{Con} (u) \rightarrow 0\tag{7}
\]

\[
0 \rightarrow \mathrm{C} \xrightarrow {\pi} \operatorname{Con} (u) \xrightarrow {\delta} \mathrm{C} ^ {\prime} (- 1) \rightarrow 0
\]

are exact.

\begin{enumerate}
\def\labelenumi{\alph{enumi})}
\setcounter{enumi}{2}
\tightlist
\item
  The morphisms \(\alpha : \mathbf{C} \to \operatorname{Cyl}(u)\) and \(\beta : \operatorname{Cyl}(u) \to \mathbf{C}\) are mutually inverse homotopisms up to homotopy.
\end{enumerate}

Assertion a) is equivalent to the formulas

\[
\tilde {u} \circ d ^ {\prime} = \overline {{\mathbf {D}}} \circ \tilde {u}, \alpha \circ d = \overline {{\mathbf {D}}} \circ \alpha , \beta \circ \overline {{\mathbf {D}}} = d \circ \beta , \pi \circ d = \mathbf {D} \circ \pi ,
\]

\[
\tilde {\pi} \circ \overline {{\mathbf {D}}} = \mathbf {D} \circ \tilde {\pi}, \delta \circ \mathbf {D} = - d ^ {\prime} \circ \delta , u = \beta \circ \tilde {u}, \pi = \tilde {\pi} \circ \alpha , \beta \circ \alpha = 1 _ {\mathrm{c}}
\]

which are verified by immediate calculations. Assertion b) is trivial. Let us prove c); on the one hand one has \(\beta \circ \alpha = 1_{c}\) ; on the other hand if \(\sigma: \text{Cyl}(u) \to \text{Cyl}(u)\) is the graded A-linear map of degree 1 such that \(\sigma(x', y', x) = (0, x', 0)\) , one verifies immediately that

\[
\overline {{{D}}} \circ \sigma + \sigma \circ \overline {{{D}}} + \alpha \circ \beta = 1 _ {\mathrm{Cyl} (u)},
\]

whence c).

One can summarize Prop. 7 by the following commutative diagram, where the rows are exact, and where the vertical arrows are homotopisms:

\[
\begin{array}{c c c c c} 0 & \longrightarrow & C & \stackrel {{\pi}} {{\longrightarrow}} & \operatorname{Con} (u) \stackrel {{\delta}} {{\longrightarrow}} C ^ {\prime} (- 1) \longrightarrow 0 \\ & & \Big \downarrow^ {\alpha} & & \Big \downarrow^ {1} \\ 0 & \longrightarrow & C ^ {\prime} \stackrel {{\bar {u}}} {{\longrightarrow}} & \operatorname{Cyl} (u) \stackrel {{\tilde {\pi}}} {{\longrightarrow}} & \operatorname{Con} (u) \longrightarrow 0 \\ & & \Big \downarrow^ {1} & & \Big \downarrow^ {\beta} \\ & & C ^ {\prime} \stackrel {{u}} {{\longrightarrow}} & C \end{array}
\]

COROLLARY. --- For every morphism of complexes \(u: \mathbf{C}' \to \mathbf{C}\) , there exists an injective morphism of complexes \(\tilde{u}: \mathbf{C}' \to \mathbf{C}_1\) and a homotopism \(\beta: \mathbf{C}_1 \to \mathbf{C}\) such that \(u = \beta \circ \tilde{u}\) .

Lemma 3. --- a) The connecting homomorphism

\[
\partial_ {n + 1} (\pi , \delta): \mathrm{H} _ {n} (\mathrm{C} ^ {\prime}) \rightarrow \mathrm{H} _ {n} (\mathrm{C})
\]

relative to the exact sequence (7) is equal to - H \(_{n}\) (u).

\begin{enumerate}
\def\labelenumi{\alph{enumi})}
\setcounter{enumi}{1}
\tightlist
\item
  The connecting homomorphism
\end{enumerate}

\[
\partial_ {n} (\tilde {u}, \tilde {\pi}): \mathrm{H} _ {n} (\operatorname{Con} (u)) \rightarrow \mathrm{H} _ {n - 1} (\mathrm{C} ^ {\prime})
\]

relative to the exact sequence (6) is equal to H \(_{n}\) (δ).

Let \(x' \in Z_n(\mathbf{C}')\) ; since \(x' = \delta(x', 0)\) and

\[
- \mathrm{D} (x ^ {\prime}, 0) = (d ^ {\prime} x ^ {\prime}, u (x ^ {\prime})) = (0, u (x ^ {\prime})) = \pi (u (x ^ {\prime})),
\]

\(\partial (\pi ,\delta)\) maps by definition the class of \(x^{\prime}\) in \(\mathbf{H}_n(\mathbf{C}')\) onto the class of \(-u(x^{\prime})\) in \(\mathbf{H}_n(\mathbf{C})\) , whence a).

Let \((y', x) \in \operatorname{Con}_n(u)\) such that \(\mathbf{D}(y', x) = 0\) ; we then have \((-d'y', dx - u(y')) = 0\) . Since \((y', x) = \tilde{\pi}(0, y', x)\) and

\[
\overline {{{\mathrm{D}}}} (0, y ^ {\prime}, x) = \left(y ^ {\prime}, - d ^ {\prime} y ^ {\prime}, d x - u (y ^ {\prime})\right) = \left(y ^ {\prime}, 0, 0\right) = \tilde {u} \big (\delta (y ^ {\prime}, x) \big),
\]

\(\partial (\tilde{u},\tilde{\pi})\) maps by definition the class of \((y^{\prime},x)\) in \(\mathrm{H}_n(\mathrm{Con}(u))\) onto the class of \(\delta (y^{\prime},x)\) in \(\mathrm{H}_{n - 1}(\mathrm{C}')\) , whence \(b)\) .

PROPOSITION 8. --- One has the unbounded exact sequence

\[
\dots \longrightarrow \mathrm{H} _ {n} (\mathrm{C} ^ {\prime}) \xrightarrow {\mathrm{H} _ {n} (u)} \mathrm{H} _ {n} (\mathrm{C}) \xrightarrow {\mathrm{H} _ {n} (\pi)} \mathrm{H} _ {n} (\text { Con } (u)) \xrightarrow {\mathrm{H} _ {n} (\delta)} \mathrm{H} _ {n - 1} (\mathrm{C} ^ {\prime}) \longrightarrow \dots .\tag{8}
\]

Indeed, in view of Lemma 3, a), this follows from Theorem 1 of X, p.~30, applied to the exact sequence (7).

COROLLARY. --- For u to be a homologism, it is necessary and sufficient that Con(u) be of zero homology.

Remark. --- Consider the diagram

\[
\begin{array}{c c c c c} \dots \longrightarrow \mathrm{H} _ {n} (\mathrm{C} ^ {\prime}) & \xrightarrow {\mathrm{H} _ {n} (\tilde {u})} \mathrm{H} _ {n} (\text {Cyl} (u)) & \xrightarrow {\mathrm{H} _ {n} (\tilde {\pi})} \mathrm{H} _ {n} (\text {Con} (u)) & \xrightarrow {\partial_ {n} (\tilde {u} , \tilde {\pi})} \mathrm{H} _ {n - 1} (\mathrm{C} ^ {\prime}) \longrightarrow \dots \\ \downarrow^ {1} & \xrightarrow {\mathrm{H} _ {n} (\beta)} \downarrow & \downarrow^ {1} & \downarrow^ {1} \\ \dots \longrightarrow \mathrm{H} _ {n} (\mathrm{C} ^ {\prime}) & \xrightarrow {\mathrm{H} _ {n} (u)} \mathrm{H} _ {n} (\mathrm{C}) & \xrightarrow {\mathrm{H} _ {n} (\pi)} \mathrm{H} _ {n} (\text {Con} (u)) & \xrightarrow {\mathrm{H} _ {n} (\delta)} \mathrm{H} _ {n - 1} (\mathrm{C} ^ {\prime}) \longrightarrow \dots \end{array}
\]

where the first line (resp. the second) is the exact homology sequence associated with the exact sequence (6) (resp. (7)). The maps \(H_{n}(\beta)\) are bijective (prop. 7, c)) and the diagram is commutative, since

\begin{enumerate}
\def\labelenumi{\alph{enumi})}
\tightlist
\item
  \(u = \beta \circ \tilde{u}\) (prop. 7, a)) therefore \(\mathrm{H}_n(u) = \mathrm{H}_n(\beta) \circ \mathrm{H}_n(\tilde{u})\) b) \(\mathrm{H}_{n}(\beta)=\mathrm{H}_{n}(\alpha)^{-1}\;\mathrm{and}\;\pi=\tilde{\pi}\circ\alpha\;(\mathrm{prop.~7,~a})\;\mathrm{and~c}),\;\mathrm{hence~}\mathrm{H}_{n}(\tilde{\pi})=\mathrm{H}_{n}(\pi)\circ\mathrm{H}_{n}(\beta),\) c) \(\mathrm{H}_{n}(\delta)=\partial_{n}(\tilde{u},\tilde{\pi})\) (Lemma 3, b)).
\end{enumerate}

